\subsection{关于波兰空间的补充}

引理 1. --- 设 X 是一个拓扑空间，A 是 X 的一个子集。为使 A 是贫集，必要且充分的是：存在 X 的一个开覆盖 \((\mathrm{U}_{i})_{i\in\mathrm{I}}\)，使 \(A\cap U_{i}\) 在 \(U_{i}\) 中贫（对每个 \(i\in I\)）。

设 \(\mathcal{O}\) 是 X 中使 \(A \cap U\) 贫的开子集 U 所成的集合。由两两不交的开集组成的 \(\mathcal{O}\) 的子集所成的族，按包含关系是归纳的。设 \(\mathfrak{U}\) 是一个极大元；这样一个元素的存在性由 E, III, p. 20, 定理 2 得出。

设 O 为 X 中属于 \(\mathfrak{U}\) 的诸开集之并。对每个开集 \(U \in \mathfrak{U}\)，设 \((B_{U,n})\) 是 U 的无处稠密子集的一个序列，其并为 \(A \cap U\)。对每个整数 n，设 \(B_n = \bigcup_{U \in \mathfrak{U}} B_{U,n}\)。

当 U 取遍诸 U 时，它相对于 O 是无处稠密的（TG, IX, p. 52, 命题 1）。由 TG, IX, p. 53, 命题 2，\(B_{n}\) 是 X 的一个无处稠密子集，因为 O 在 X 中开。于是，等于诸 \(B_{n}\) 之并的集合 A ∩ O 是 X 的一个贫子集。

设 F 为 O 的补集。它是 X 的闭子集；我们来证明它的内部为空。若不然，设 x 是 F 的一个内点。依假设，存在 x 的一个开邻域 V——可设它包含于 F——使 A ∩ V 在 X 中贫。于是 V 与 X 中属于 U 的诸开集不交，且 U ∪ \{V\} 是由属于 O 的两两不交的开集组成的一个集合，这与 U 的极大性矛盾。关系式

\[
\mathrm{A} = (\mathrm{A} \cap \mathrm{F}) \cup (\mathrm{A} \cap \mathrm{O})
\]

于是蕴含 A 贫，此即所要证明的。

设 X 是一个拓扑空间。

回忆（TG, IX, p. 69）：称 X 的一个子集 A 是可达的，如果存在 X 的一个开集 U，使 \(U \cap \complement A\) 与 \(A \cap \complement U\) 都在 X 中贫。X 的可达子集的全体是一个 \(\sigma\)-代数，它包含博雷尔 \(\sigma\)-代数（TG, IX, p. 69, 引理 8 及其证明）。贫子集是可达的。

对 X 的任一子集 A，设 D(A) 为由满足下列条件的点 \(x \in X\) 组成的集合：对 x 的每个邻域 V，A∩V 都不贫。我们还置 \(\mathrm{D}^{*}(\mathrm{A}) = \mathrm{A} \cup \mathrm{D}(\mathrm{A})\)。

引理 2. --- 设 X 是一个拓扑空间，A 是 X 的一个子集。

\begin{enumerate}
\def\labelenumi{\alph{enumi})}
\item
  集合 D(A) 在 X 中闭；其补集是 X 中使 A ∩ U 为贫集的最大开集 U。
\item
  为使 A 贫，必要且充分的是 D(A) 为空。
\item
  集合 \(\mathrm{D}^{*}(\mathrm{A})\) 是可达的。
\item
  对 X 的每个包含 A 的可达集 B，\(D^{*}(A) \cap \complement B\) 贫。
\end{enumerate}

设 U 为 X 中使 \(A \cap V\) 贫的诸开集 V 之并。为使一点 x 属于 U，必要且充分的是它有一个邻域 V 使 \(A \cap V\) 贫。于是 \(U = \complement D(A)\)，因此 D(A) 是 X 的闭子集。按构造，子空间 U 有一个由与 A 的交为贫的开集组成的覆盖。把引理 1 应用于拓扑空间 U 与子集 \(A \cap U\)，可知 \(A \cap U\) 在 U 中贫，因而也在 X 中贫。这就证明了断言 a)，因为按构造，X 中使 \(A \cap V\) 贫的每个开集 V 都包含于 U。

断言 b) 由此立得。

\begin{enumerate}
\def\labelenumi{\alph{enumi})}
\setcounter{enumi}{2}
\item
  我们有 \(D^{*}(A) = A \cup D(A) = (A \cap U) \cup D(A)\)。集合 \(D(A)\) 是闭的，因而是可达的（IV, p. 361）；贫的部分 \(A \cap U\) 也是可达的。于是 \(D^{*}(A)\) 可达（同前）。
\item
  设 B 是 X 的一个包含 A 的可达子集。其补集 \(\complement B\) 于是是 X 的可达子集（同前），因而存在 X 的一个开集 V，使 \(V \cap B\) 与 \(\complement V \cap \complement B\) 贫。由于 \(A \subset B\)，\(V \cap A\) 仍贫，故 \(V \subset U\)。由于 \(\complement U \cap \complement B\) 包含于 \(\complement V \cap \complement B\)，它也是 X 的贫子集。最后，包含关系
\end{enumerate}

\[
\mathrm{D} ^ {*} (\mathrm{A}) \cap \complement \mathrm{B} = ((\mathrm{A} \cap \mathrm{U}) \cap \complement \mathrm{B}) \cup (\complement \mathrm{U} \cap \complement \mathrm{B}) \subset (\mathrm{A} \cap \mathrm{U}) \cup (\complement \mathrm{U} \cap \complement \mathrm{B})
\]

证明了 \(D^{*}(A) \cap \complement B\) 是 X 的一个贫子集。

注 1. --- 设 G 是一个拓扑群，A 是 G 的一个贫子集。假设 A 包含 G 的一个非空开集 U，并设 y 是 U 的一点。那么 U 贫，且对每个 \(x \in G\)，\(xy^{-1}U\) 是 x 在 G 中的邻域。于是 G 的每一点都有一个贫邻域，故 G 贫（IV, p. 360, 引理 1）。

反之，若 G 不贫，则每个贫子集的内部为空，且 G 是贝尔空间。由于贝尔空间不贫，这就证明了下列断言：为使一个拓扑群是贝尔空间，必要且充分的是它不贫。

命题 7. --- 设 X 是一个豪斯多夫空间。X 的每个苏斯林子空间（TG, IX, p. 59, 定义 2）都是可达的。

设 S 是 X 的一个苏斯林子空间。按定义，存在一个具有可数基的完备度量空间 P 及一个从 P 到 S 上的连续满射 g。由 TG, IX, p. 64, 引理 3，存在度量空间 P 的一个筛 \(\mathrm{C} = (\mathrm{C}_{n}, p_{n}, \varphi_{n})_{n \in \mathbf{N}}\)。

对每个整数 \(n\) 及每个 \(c \in \mathrm{C}_n\)，记 \(\mathrm{F}_n(c)\) 为集合 \(\varphi_n(c)\) 在 \(g\) 下的像；再置 \(\mathrm{F}_n^*(c) = \mathrm{D}^*(\mathrm{F}_n(c))\)，以及

\[
\mathrm{G} _ {n} (c) = \mathrm{F} _ {n} ^ {*} (c) \cap \mathbb {C} \left(\bigcup_ {c ^ {\prime} \in p _ {n} ^ {- 1} (c)} \mathrm{F} _ {n + 1} ^ {*} (c ^ {\prime})\right).\tag{1}
\]

对每个 \(c \in C_{n}\)，\(\mathrm{F}_{n}^{*}(c)\) 是可达的（IV, p. 361, 引理 2, c)）。由于 \(C_{n+1}\) 可数，诸可达子集 \(\mathrm{F}_{n+1}^{*}(c')\)（\(c'\) 取遍 \(p_{n}^{-1}(c)\)）之并仍是 X 的一个可达子集。它包含诸 \(\mathrm{F}_{n+1}(c')\) 之并，后者等于 \(\mathrm{F}_{n}(c)\)。于是（同前， d)），\(\mathrm{G}_{n}(c)\) 是 X 的一个贫子集。因此，诸集合 \(\mathrm{G}_{n}(c)\)（对 \(n \in \mathbf{N}\) 及 \(c \in C_{n}\)）之并 G 是 X 的一个贫子集。

设 \(c_0 \in \mathrm{C}_0\)，设 \(x\) 是 \(\mathrm{F}_0^*(c_0) \cap \complement \mathrm{G}\) 的一个元素；我们来证明 \(x \in \mathrm{F}_0(c_0)\)。由于 \(x \notin \mathrm{G}_0(c_0)\) 且 \(x \in \mathrm{F}_0^*(c_0)\)，由关系式 (1)，存在 \(c_1 \in p_0^{-1}(c_0)\) 使 \(x \in \mathrm{F}_1^*(c_1)\)。由归纳法，存在一个元素 \(c = (c_n)_n\)（属于 \(\prod_n C_n\)），使得对每个整数 \(n \in \mathbf{N}\)，有 \(x \in \mathrm{F}_n^*(c_n)\) 且 \(p_n(c_{n+1}) = c_n\)。

诸子集 \(\varphi_n(c_n)\) 构成 P 中一个柯西滤子的基；由于 P 完备，这个滤子收敛于 P 的一点 \(p\)。这个滤子在 \(g\) 下的像是 X 上的一个滤子 F，以诸 \(\mathrm{F}_n(c_n)\)（\(n\in \mathbf{N}\)）组成的集合为基，且收敛于 \(g(p)\)。由于 \(\mathrm{F}_n^* (c_n)\) 包含于 \(\overline{\mathrm{F}_n(c_n)}\) 且 \(x\) 属于每个 \(\mathrm{F}_n(c_n)\)，\(x\) 是 F 的一个附着点，于是 \(x = g(p)\)，因为空间 X 是豪斯多夫的（TG, I, p. 52, 命题 1）。点 \(p\) 属于 \(\overline{\varphi_1(c_1)}\)，因而属于 \(\varphi_0(c_0)\)。于是，\(x\in \mathrm{F}_0(c_0)\)。

于是，集合 \(\mathrm{F}_{0}^{*}(c_{0}) - \mathrm{F}_{0}(c_{0})\) 包含于 G。它因而是 X 的一个贫子集，从而也是可达子集。由于 \(\mathrm{F}_{0}(c_{0}) = \mathrm{F}_{0}^{*}(c_{0}) - (\mathrm{F}_{0}^{*}(c_{0}) - \mathrm{F}_{0}(c_{0}))\) 且 \(\mathrm{F}_{0}^{*}(c_{0})\) 可达，故 \(\mathrm{F}_{0}(c_{0})\) 可达，因为 X 的可达子集构成一个 \(\sigma\)-代数。由于 \(C_{0}\) 可数，作为诸 \(\mathrm{F}_{0}(c_{0})\)（\(c_{0} \in C_{0}\)）之并的集合 S 仍是可达的。

